\SourceLeaf{43}\begin{center}
{\Large SECTION THREE}\par\bigskip
{\small ON}\par\bigskip
{\Large RESIDUES OF POWERS.}
\end{center}

\begin{center}\rule{0.09\linewidth}{0.3pt}\end{center}

\begin{center}\small\emph{The residues of the terms of a geometric progression beginning with unity form a periodic sequence.}\end{center}

\Article{45}

\textsc{Theorem.} \emph{In every geometric progression $1,a,aa,a^3$ etc., besides the first term $1$, there is another term $a^t$ congruent to unity modulo a number $p$ coprime to $a$, with exponent $t<p$.}

\emph{Proof.} Since the modulus $p$ is coprime to $a$, and therefore to every power of $a$, no term of the progression is $\equiv0\ (\mathrm{mod.}p)$. Each is congruent to one of the numbers $1,2,3....p-1$. There are $p-1$ such numbers, so if more than $p-1$ terms of the progression are considered, they cannot all have distinct least residues. Thus at least two of the terms $1,a,aa,a^3....a^{p-1}$ are congruent. Let $a^m\equiv a^n$, with $m>n$. Dividing by $a^n$ gives $a^{m-n}\equiv1$ (Article 22), where $m-n<p$ and $>0$.\par
\noindent\mbox{\emph{Q. E. D.}}

\emph{Example.} In the progression $2,4,8$ etc., the first term congruent to unity modulo $13$ is $2^{12}=4096$. But modulo $23$, the same progression gives $2^{11}=2048\equiv1$. Similarly, the sixth power of $5$, namely $15625$, is congruent to unity modulo $7$, and the fifth power, $3125$, modulo $11$. So in some cases a power with exponent less than $p-1$ is congruent to unity; in others we must go as far as the power with exponent $p-1$.

\SourceLeaf{44}\Article{46}

If the progression is continued beyond a term congruent to unity, the residues from the beginning occur again. If $a^t\equiv1$, then $a^{t+1}\equiv a,a^{t+2}\equiv aa$ etc., until the term $a^{2t}$ is reached. Its least residue is again $\equiv1$, and it starts the \emph{period} of residues once more. Thus there is a period containing $t$ residues, which always repeats from the beginning as soon as it ends. No residues other than those in this period occur anywhere in the progression. In general, $a^{mt}\equiv1$ and $a^{mt+n}\equiv a^n$, which in our notation is expressed as follows:

\[\text{\emph{If }}r\equiv\rho\ (mod.\,t),\quad\text{\emph{then}}\quad a^r\equiv a^\rho\ (mod.\,p).\]

\Article{47}

This theorem gives a shortcut for finding residues of powers, however large their exponents, once a power congruent to unity is known. For \emph{example}, to find the residue obtained by dividing $3^{1000}$ by $13$, we have, since $3^3\equiv1\ (\mathrm{mod.}\ 13)$, $t\equiv3$. Since $1000\equiv1\ (\mathrm{mod.}\ 3)$, it follows that $3^{1000}\equiv3\ (\mathrm{mod.}\ 13)$.

\Article{48}

If $a^t$ is the \emph{lowest} power congruent to unity, apart from $a^0=1$, which we leave out of consideration here, the $t$ terms making up the period of residues are all distinct, as the proof of Article 45 readily shows. We can then reverse the proposition of Article 46: if $a^m\equiv a^n\ (\mathrm{mod.}p)$, then $m\equiv n\ (\mathrm{mod.}t)$. For if $m,n$ were incongruent modulo $t$, their least residues $\mu,\nu$ would differ. But $a^\mu\equiv a^m,a^\nu\equiv a^n$, so $a^\mu\equiv a^\nu$: \emph{that is}, the powers below $a^t$ would not all be incongruent, contrary to the hypothesis.

So, if $a^k\equiv1\ (\mathrm{mod.}p)$, then $k\equiv0\ (\mathrm{mod.}t)$, \emph{that is}, $k$ is divisible by $t$.

So far we have considered any moduli, provided they are coprime to $a$. We will now consider prime moduli separately, and later build a more general investigation on this foundation.

\SourceLeaf{45}\begin{center}\small\emph{First we consider moduli that are prime numbers.}\end{center}

\Article{49}

\textsc{Theorem.} \emph{If $p$ is prime and does not divide $a$, and $a^t$ is the lowest power of $a$ congruent to unity modulo $p$, then the exponent $t$ is either $=p-1$ or a divisor of that number.}

Compare the examples in Article 45.

\emph{Proof.} We have already shown that $t$ is either $=p-1$ or $<p-1$. It remains to prove that in the latter case $t$ always divides $p-1$.

I. Take the least positive residues of all the terms $1,a,aa....a^{t-1}$ and denote them by $\alpha,\alpha',\alpha''$ etc., so that $\alpha=1,\alpha'\equiv a,\alpha''\equiv aa$ etc. They are clearly all distinct: if two terms $a^m,a^n$ gave the same residue, then, assuming $m>n$, we would have $a^{m-n}\equiv1$ with $m-n<t$, \mbox{\emph{Q. E. A.}}, since no power below $a^t$ is congruent to unity (\emph{by hypothesis}). All the numbers $\alpha,\alpha',\alpha''$ etc. are in the sequence $1,2,3...p-1$, but do not exhaust it, since $t<p-1$. Denote the collection of all $\alpha,\alpha',\alpha''$ etc. by $(A)$. Thus $(A)$ contains $t$ terms.

II. Choose any number $\beta$ from $1,2,3...p-1$ that is absent from $(A)$. Multiply $\beta$ by all of $\alpha,\alpha',\alpha''$ etc., and let the resulting least residues be $\beta,\beta',\beta''$ etc. Their number is again $t$. These residues are distinct both from each other and from all of $\alpha,\alpha',\alpha''$ etc. If the \emph{first} assertion were false, we would have $\beta a^m\equiv\beta a^n$, and division by $\beta$ would give $a^m\equiv a^n$, contrary to what we just proved. If the \emph{second} were false, we would have $\beta a^m\equiv a^n$. When $m<n$, this gives $\beta\equiv a^{n-m}$, \emph{that is}, $\beta$ would be congruent to one of $\alpha,\alpha',\alpha''$ etc., contrary to the hypothesis. When $m>n$, multiplication by $a^{t-m}$ gives $\beta a^t\equiv a^{t+n-m}$, or, since $a^t\equiv1$, $\beta\equiv a^{t+n-m}$, the same contradiction. Denote the collection of all $\beta,\beta',\beta''$ etc., whose number is $=t$, by $(B)$. We now have $2t$ numbers from $1,2,3...p-1$. If $(A)$ and $(B)$ include all these numbers, then $\frac{p-1}{2}=t$, and the theorem is proved.

III. If some numbers are still missing, let $\gamma$ be one of them. Multiply all of $\alpha,\alpha',\alpha''$ etc. by it, and let the least residues of the products be $\gamma,\gamma',\gamma''$ etc. Denote their collection by $(C$. Then $(C)$ contains $t$ numbers from $1,2,3...p-1$, all distinct from each other and from the numbers in $(A)$ and $(B)$.
\SourceLeaf{46}
The first two assertions are proved as in II; the third as follows. If $\gamma a^m\equiv\beta a^n$, then $\gamma\equiv\beta a^{n-m}$ or $\equiv\beta a^{t+n-m}$, according as $m<n$ or $>n$. In either case $\gamma$ would be congruent to a member of $(B)$, contrary to the hypothesis. We thus have $3t$ numbers from $1,2,3...p-1$. If none are missing, $t=\frac{p-1}{3}$, and the theorem is proved.

IV. If some numbers are still missing, proceed in the same way to a fourth collection $(D)$, and so on. Since there are only finitely many numbers $1,2,3....p-1$, they must eventually be exhausted. Their number is therefore a multiple of $t$, so $t$ divides $p-1$. \mbox{\emph{Q. E. D.}}

\Needspace{6\baselineskip}
\begin{center}\small\emph{Fermat's theorem.}\end{center}

\Article{50}

Since $\frac{p-1}{t}$ is an integer, raising both sides of the congruence $a^t\equiv1$ to the power with exponent $\frac{p-1}{t}$ gives $a^{p-1}\equiv1$. \emph{In other words, $a^{p-1}-1$ is always divisible by $p$ when $p$ is prime and does not divide $a$.}

This theorem deserves close attention for both its elegance and its exceptional usefulness. It is usually called \emph{Fermat's theorem}, after its discoverer; see \emph{Fermatii Opera Mathem. Tolosae} 1679, \emph{fol. p.} 163. He gave no proof, though he said he had one. The distinguished Euler was the first to publish a proof, in his paper \emph{Theorematum quorundam ad numeros primos spectantium demonstratio, Comm. Acad. Petrop. T.} VIII\begingroup\renewcommand{\thefootnote}{*)}\footnote{In an earlier paper this great man had not yet reached his goal: \emph{Comm. Petr. T.} VI, \emph{p.} 106. --- In the famous controversy between Maupertuis and Konig, which began over the principle of least action but soon turned to unrelated matters, K\"onig said he possessed a manuscript in Leibniz's own hand containing a proof of this theorem entirely agreeing with Euler's: \emph{Appel au public. p.} $106$. Although we do not wish to deny the credibility of this testimony, Leibniz certainly never published his discovery. Compare \emph{Hist. de l'Ac. de Prusse, A.} 1750, \emph{p.} 530.}\endgroup. It rests on expanding $(a+1)^p$. The form of the coefficients readily shows that $(a+1)^p-a^p-1$ is always divisible by $p$. So $(a+1)^p-(a+1)$ is divisible by $p$ whenever $a^p-a$ is divisible by $p$. Since $1^p-1$ is always divisible by $p$, so is $2^p-2$, and therefore $3^p-3$ etc., and in general $a^p-a$. If $p$ does not divide $a$, then $a^{p-1}-1$ is also divisible by $p$. This is enough to show the character of the method. The distinguished Lambert gave a similar proof in \emph{Actis Erudit.} 1769,
\SourceLeaf{47}
\emph{p.} 109. But expanding a binomial power seemed rather foreign to number theory, so the distinguished Euler sought another proof. It appears in \emph{Comment. nov. Petr. T.} VII, \emph{p.} 70, and entirely agrees with the proof we gave in the preceding article. We will meet some further proofs later. Here we may add one based on principles similar to those of Euler's first proof. The following proposition, of which our theorem is only a special case, will also be used in other investigations below.

\Article{51}

\emph{The polynomial $a+b+c+\mathit{etc.}$, raised to the power with exponent $p$, is, modulo $p$,}

\[\equiv a^p+b^p+c^p+\mathit{etc.}\]

\noindent\emph{provided that $p$ is prime.}

\emph{Proof.} The power with exponent $p$ of the polynomial $a+b+c+\text{etc.}$ consists of terms of the form $\varkappa a^\alpha b^\beta c^\gamma$ etc., where $\alpha+\beta+\gamma\text{ etc.}=p$, and $\varkappa$ denotes the number of permutations of $p$ objects, of which $\alpha,\beta,\gamma$ etc. are respectively $=a,b,c$ etc. In Article 41 we showed that this number is always divisible by $p$, unless all the objects are equal, \emph{that is}, unless one of the numbers $\alpha,\beta,\gamma$ etc. is $=p$ and the others $=0$. So all terms of $(a+b+c+\text{etc.})^p$, except $a^p,b^p,c^p$ etc., are divisible by $p$. When considering congruence modulo $p$, we can therefore discard them, giving

\[(a+b+c+\mathrm{etc.})^p\equiv a^p+b^p+c^p+\mathrm{etc.}\quad\text{\emph{Q. E. D.}}\]

If all the quantities $a,b,c$ etc. are now set $=1$, and their number $=k$, this gives $k^p\equiv k$, as in the preceding article.

\Needspace{6\baselineskip}
\begin{center}\emph{How many numbers give periods whose number of terms is a given divisor of $p-1$.}\end{center}

\Article{52}

Only divisors of $p-1$ can be the exponents of the lowest powers to which numbers are raised to become congruent to unity. So we ask whether all divisors of $p-1$ can occur, and how many numbers belong to each exponent when all numbers not divisible by $p$ are classified according to the exponent of their lowest power congruent to unity. We should first note that it is enough to consider the positive numbers
\SourceLeaf{48}
from $1$ to $p-1$. Congruent numbers clearly have to be raised to the same lowest power to become congruent to unity, so any number belongs to the same exponent as its least positive residue. We must therefore find how, in this respect, the numbers $1,2,3\GaussFourDots p-1$ are distributed among the individual factors of $p-1$. For brevity, if $d$ is a divisor of $p-1$, including $1$ and $p-1$, let $\psi d$ denote the number of positive numbers smaller than $p$ whose lowest power congruent to unity has exponent $d$.

\Article{53}

To make the investigation easier to follow, we give an example. For $p=19$, the numbers $1,2,3\GaussFourDots18$ are distributed among the divisors of $18$ as follows:

\[
\begin{array}{r|l}
1 & \phantom{1}1.\\
2 & 18.\\
3 & \phantom{1}7,\ 11.\\
6 & \phantom{1}8,\ 12.\\
9 & \phantom{1}4,\ 5.\ 6,\ 9,\ 16,\ 17.\\
18 & \phantom{1}2,\ 3,\ 10,\ 13.\ 14,\ 15.
\end{array}
\]

Thus, in this case, $\psi1=1$, $\psi2=1$, $\psi3=2$, $\psi6=2$, $\psi9=6$, $\psi18=6$. A little attention shows that as many numbers belong to each exponent as there are numbers no greater than it and coprime to it. So, at least in this case, using the notation of Article 39, $\psi d=\phi d$. We prove that this observation holds generally as follows.

I. Suppose we have a number $a$ \emph{belonging} to exponent $d$, \emph{that is}, whose power with exponent $d$ is congruent to unity, while all lower powers are incongruent to it. All its powers $aa,a^3,a^4\GaussFourDots a^d$, or their least residues, also have the first property: their powers with exponent $d$ are congruent to unity. We can also say that the least residues of $a,aa,a^3\GaussFourDots a^d$, all distinct, are roots of the congruence $x^d\equiv1$. Since this congruence cannot have more than $d$ distinct roots, no numbers other than the least residues of $a,aa,a^3\GaussFourDots a^d$ exist between $1$ and $p-1$, inclusive, whose powers with exponent\SourceLeaf{49} $d$ congruent to unity. Thus all numbers belonging to exponent $d$ are among the least residues of $a,aa,a^3\GaussFourDots a^d$. We determine which ones they are, and how many, as follows. If $k$ is coprime to $d$, no power of $a^k$ with exponent $<d$ is congruent to unity. For let $\frac{1}{k}\,(\operatorname{mod.}d)\equiv m$ (see Article 31); then $a^{km}\equiv a$. If the power with exponent $e$ of $a^k$ were congruent to unity with $e<d$, we would also have $a^{kme}\equiv1$, and therefore $a^e\equiv1$, contrary to the hypothesis. So the least residue of $a^k$ clearly belongs to exponent $d$. If, however, $k$ has a common divisor $\delta$ with $d$, the least residue of $a^k$ does not belong to exponent $d$, since its power with exponent ${\frac{d}{\delta}}$ is already congruent to unity. For $\frac{kd}{\delta}$ is divisible by $d$, or $\equiv0\,(\operatorname{mod.}d)$, and so $a^{\frac{kd}{\delta}}\equiv1$. Thus as many numbers belong to exponent $d$ as there are numbers among $1,2,3\GaussFourDots d$ coprime to $d$. But remember that this conclusion assumes we already have one number $a$ belonging to exponent $d$. It remains uncertain whether there could be an exponent to which no number belongs at all. The conclusion is therefore limited to saying that $\psi d$ is either $=0$ or $=\phi d$.

\Article{54}

II. Let all the divisors of $p-1$ be $d,d',d''$ etc. Since all numbers $1,2,3\GaussFourDots p-1$ are distributed among these, we have

\[\psi d+\psi d'+\psi d''+\text{etc.}=p-1\]

But in Article 40 we proved that

\[\phi d+\phi d'+\phi d''+\text{etc.}=p-1\]

The preceding article shows that $\psi d$ is equal to or smaller than $\phi d$, and cannot be larger. The same holds for $\psi d'$ and $\phi d'$, etc. If any of $\psi d,\psi d',\psi d''$ etc. were smaller than the corresponding term among $\phi d,\phi d',\phi d''$, or if several were, the two sums could not be equal. We finally conclude that $\psi d$ \emph{is always equal to $\phi d$}, and therefore does not depend on the size of $p-1$.

\Article{55}

One special case of the preceding\SourceLeaf{50} proposition deserves particular attention: \emph{there are always numbers for which no power with exponent below $p-1$ is congruent to unity}. There are as many of them between $1$ and $p-1$ as there are numbers below $p-1$ coprime to $p-1$. The proof of this theorem is far from as obvious as it may seem at first. Because of its importance, we may add another proof somewhat different from the preceding one: different methods often do much to clarify more difficult matters. Resolve $p-1$ into its prime factors, writing $p-1=a^\alpha b^\beta c^\gamma$ etc., where $a,b,c$ etc. are distinct primes. We will complete the proof by establishing the following:

I. We can always find a number $A$, or several, belonging to exponent $a^\alpha$, and similarly numbers $B,C$ etc. belonging to exponents $b^\beta,c^\gamma$ etc., respectively.

II. The product of all the numbers $A,B,C$ etc., or its least residue, belongs to exponent $p-1$. We prove these as follows.

I. Let $g$ be one of the numbers $1,2,3\ldots p-1$ \emph{not} satisfying the congruence $x^{\frac{p-1}{a}}\equiv1\,(\operatorname{mod.}p)$, since these numbers cannot all satisfy a congruence of degree $<p-1$. I claim that, if the power with exponent ${\frac{p-1}{a^\alpha}}$ of $g$ is set $\equiv h$, this number, or its least residue, belongs to exponent $a^\alpha$.

For the power with exponent $a^{\alpha}$ of $h$ is clearly congruent to the power with exponent $p-1$ of $g$, \emph{that is}, to unity. But the power with exponent $a^{\alpha-1}$ of $h$ is congruent to the power with exponent ${\frac{p-1}{a}}$ of $g$, \emph{that is}, it is incongruent to unity. Still less can the powers with exponents $a^{\alpha-2},\ a^{\alpha-3}$ etc. of $h$ be congruent to unity. Yet the exponent of the lowest power of $h$ congruent to unity, or the exponent to which $h$ belongs, must divide $a^\alpha$ (Article 48). Since $a^\alpha$ is divisible only by itself and lower powers of $a$, $a^\alpha$ must be the exponent to which $h$ belongs. \mbox{\emph{Q. E. D.}} The same method proves the existence of numbers belonging to exponents $b^\beta,c^\gamma$ etc.

II. Suppose the product of all $A,B,C$ etc. belongs not to exponent $p-1$ but to a smaller exponent $t$. Then $t$ divides $p-1$ (Article 48), so $\frac{p-1}{t}$ is an integer greater than unity. This quotient is clearly either one of the primes $a,b,c$ etc. or at least divisible by one of them (Article 17), for \emph{example} by $a$; the reasoning is similar for the others.

\SourceLeaf{51}
So $t$ divides $\frac{p-1}{a}$, and the product $ABC$ etc. is also congruent to unity when raised to the power with exponent ${\frac{p-1}{a}}$ (Article 46). Each of $B,C$ etc., excluding $A$, is clearly congruent to unity when raised to the power with exponent ${\frac{p-1}{a}}$, since the exponents $b^\beta,c^\gamma$ etc. to which they belong divide $\frac{p-1}{a}$. Thus
\[A^{\frac{p-1}{a}}B^{\frac{p-1}{a}}C^{\frac{p-1}{a}}\text{ etc.}\equiv A^{\frac{p-1}{a}}\equiv1.\]
So the exponent to which $A$ belongs must divide $\frac{p-1}{a}$ (Article 48), \emph{that is}, $\frac{p-1}{a^{\alpha+1}}$ must be an integer. But $\frac{p-1}{a^{\alpha+1}}=\frac{b^\beta c^\gamma\text{ etc.}}{a}$ cannot be an integer (Article 15). We must therefore conclude that our assumption is impossible, \emph{that is}, the product $ABC$ etc. really does belong to exponent $p-1$. \mbox{\emph{Q. E. D.}}

The second proof seems somewhat longer than the first; the first, on the other hand, is less direct than the second.

\Article{56}
This theorem is a striking example of the care often needed in number theory to avoid taking things for certain when they are not. The celebrated Lambert mentions this proposition in the paper cited above, \emph{Acta Erudit.} 1769, \emph{p.} 127, but does not even touch on the need for a proof. No one attempted a proof except the great Euler: \emph{Comment. nov. Ac. Petrop. T.} XVIII, for 1773, \emph{Demonstrationes circa residua ex divisione potestatum per numeros primos resultantia}, p. 85 onwards. See especially Article 37, where he discusses the need for a proof at greater length. But the proof given by this most perceptive man has two gaps. First, in Article 31 and the following articles, he tacitly assumes that the congruence $x^n\equiv1$, expressing his reasoning in our notation, really has $n$ distinct roots, though previously he had shown only that it cannot have \emph{more}. Second, he derived the formula in Article 34 by induction alone.

\Needspace{6\baselineskip}
\begin{center}\emph{Primitive roots, bases, indices.}\end{center}

\Article{57}
Following the distinguished Euler, we call numbers belonging to exponent $p-1$ \emph{primitive roots}. If $a$ is a primitive root, the least residues of the powers $a,aa,a^3\ldots a^{p-1}$
\SourceLeaf{52}
are all distinct. So they must include all the numbers $1,2,3,\ldots p-1$, since there are just as many of these as of the least residues. In other words, every number not divisible by $p$ is congruent to a power of $a$. This remarkable property is very useful, and greatly eases the arithmetic of congruences, much as logarithms ease ordinary arithmetic. We choose any primitive root $a$ as a \emph{base}, to which we refer all numbers not divisible by $p$. If $a^e\equiv b\,(\operatorname{mod.}p)$, we call $e$ the \emph{index} of $b$. For \emph{example}, for modulus $19$, taking the primitive root $2$ as base gives the following correspondence:

\begin{center}
\begingroup\small\setlength{\tabcolsep}{1.5pt}
\begin{tabular}{@{}l@{\quad}*{18}{r}@{}}
numbers & 1. & 2. & 3. & 4. & 5. & 6. & 7. & 8. & 9. & 10. & 11. & 12. & 13. & 14. & 15. & 16. & 17. & 18.\\
indices & 0. & 1. & 13. & 2. & 16. & 14. & 6. & 3. & 8. & 17. & 12. & 15. & 5. & 7. & 11. & 4. & 10. & 9.
\end{tabular}
\endgroup
\end{center}

For a fixed base, each number clearly has several indices, but they are all congruent modulo $p-1$. Whenever we discuss indices, those congruent modulo $p-1$ will therefore be regarded as equivalent, just as numbers themselves are regarded as equivalent when congruent modulo $p$.

\Needspace{6\baselineskip}
\begin{center}\emph{Calculation with indices.}\end{center}

\Article{58}
The theorems on indices are entirely analogous to those on logarithms.

\emph{The index of a product of any number of factors is congruent to the sum of the indices of the individual factors modulo $p-1$.}

\emph{The index of a power of a number is congruent to the product of the index of the given number and the exponent of the power, modulo $p-1$.}

We omit the proofs because they are easy.

So, in constructing a table from which the indices of all numbers can be found for different moduli, we can omit both numbers greater than the modulus and composite numbers. A specimen is appended at the end of this work, \emph{Table} I. Its first vertical column lists the primes and prime powers from $3$ to $97$, to be taken as moduli, with the chosen base beside each. The indices of successive primes follow, separated into groups of five by a small
gap\SourceLeaf{53}, with the primes arranged above in the same order. Thus the index of a given prime for a given modulus can be found easily and reliably.

For \emph{example}, if $p=67$, the index of $60$, with $12$ as base, is
\[\equiv2\operatorname{Ind.}2+\operatorname{Ind.}3+\operatorname{Ind.}5\ (\operatorname{mod.}66)\equiv58+9+39\equiv40\]

\Article{59}
\emph{The index of any value of the expression $\frac{a}{b}(\mathit{mod.}p.),$ (Article 31) is congruent modulo $p-1$ to the difference between the indices of the numerator $a$ and denominator $b$, provided that the numbers $a,b$ are not divisible by $p$.}

For let any such value be $c$. Then $bc\equiv a\,(\operatorname{mod.}p)$, so
\[\operatorname{Ind.}b+\operatorname{Ind.}c\equiv\operatorname{Ind.}a\,(\operatorname{mod.}p-1)\]
and therefore
\[\operatorname{Ind.}c\equiv\operatorname{Ind.}a-\operatorname{Ind.}b\]

If we have a table giving the index of each number for any prime modulus, and another giving the number belonging to a given index, every congruence of the first degree can be solved very easily, since all can be reduced to congruences with prime moduli (Article 30).

For \emph{example}, given the congruence
\[29x+7\equiv0\,(\operatorname{mod.}47)\quad\text{then}\quad x\equiv\frac{-7}{29}(\operatorname{mod.}47)\]
we obtain
\[\operatorname{Ind.}x\equiv\operatorname{Ind.}-7-\operatorname{Ind.}29\equiv\operatorname{Ind.}40-\operatorname{Ind.}29\equiv15-43\equiv18(\operatorname{mod.}46)\]
The number with index $18$ is $3$, so $x\equiv3\,(\operatorname{mod.}47)$. --- We have not supplied the second table, but another can take its place, as we will show in Section VI.

\Needspace{6\baselineskip}
\begin{center}\emph{Roots of the congruence $x^n\equiv A$.}\end{center}

\Article{60}
As in Article 31, where we introduced notation for roots of congruences of the first degree, we will now also use a symbol for roots of pure congruences of higher degrees. Just as ${}^{n}\!\surd A$ means a root of the equation $x^n=A$, adding a modulus gives ${}^{n}\!\surd A\,(\operatorname{mod.}p)$ to denote any root of the congruence $x^n\equiv A\,(\operatorname{mod.}p)$. We say that the expression ${}^{n}\!\surd A\,(\operatorname{mod.}p)$ has as many values as it has values incongruent modulo $p$, since all values congruent modulo $p$
\SourceLeaf{54}
are to be regarded as equivalent (Article 26). If $A,B$ are congruent modulo $p$, the expressions ${}^{n}\!\surd A,{}^{n}\!\surd B\,(\operatorname{mod.}p)$ are clearly equivalent.

If we put ${}^{n}\!\surd A\equiv x\,(\operatorname{mod.}p)$, then $n\operatorname{Ind.}x\equiv\operatorname{Ind.}A\,(\operatorname{mod.}p-1)$. The rules of the preceding section give the values of $\operatorname{Ind.}x$ from this congruence, and from them the corresponding values of $x$. Clearly $x$ has as many values as the congruence $n\operatorname{Ind.}x\equiv\operatorname{Ind.}A\,(\operatorname{mod.}p-1)$ has roots. Thus ${}^{n}\!\surd A$ has just one value when $n$ is coprime to $p-1$. If, however, the numbers $n,p-1$ have greatest common divisor $\delta$, then $\operatorname{Ind.}x$ has $\delta$ values incongruent modulo $p-1$, and ${}^{n}\!\surd A$ has just as many values incongruent modulo $p$, provided that $\operatorname{Ind.}A$ is divisible by $\delta$. If this condition fails, ${}^{n}\!\surd A$ has no real value.

\emph{Example.} Find the values of ${}^{15}\!\surd11\,(\operatorname{mod.}19)$. We must solve the congruence $15\operatorname{Ind.}x\equiv\operatorname{Ind.}11\equiv6\,(\operatorname{mod.}18)$, giving the three values $\operatorname{Ind.}x\equiv4,10,16\,(\operatorname{mod.}18)$. The corresponding values of $x$ are $6,9,4$.
